\section{Antecedentes}
\label{sec:antecedentes}

\subsection{Revisión de literatura}

Los avances recientes en ciencia de datos, estadística avanzada y aprendizaje automático han impulsado una evolución significativa en los métodos para predecir fallas, gestionar interrupciones y estimar la calidad del servicio en redes de distribución eléctrica.

En primer lugar, varios estudios han explorado la predicción de confiabilidad a partir de modelos híbridos que integran técnicas estadísticas y de inteligencia artificial. \textcite{ref11} desarrollan un modelo integral de pronóstico de confiabilidad que combina redes neuronales artificiales, análisis de relaciones grises y modelos de fallas programadas, demostrando alta aplicabilidad en sistemas en expansión. De manera complementaria, \textcite{ref12} proponen un ensamble de redes neuronales profundas que mejora sustancialmente la estimación de interrupciones causadas por viento y descargas eléctricas, superando modelos estadísticos y de comité previos.

En el ámbito de eventos meteorológicos severos, \textcite{ref13} introducen uno de los primeros modelos operativos basados en Machine Learning para la planificación de tormentas, combinando datos climáticos, propiedades físicas de la red e información histórica de daños. En la misma línea, \textcite{ref14} emplean Graph Neural Networks para procesar mediciones meteorológicas distribuidas en un marco de señales sobre grafos, logrando errores de predicción inferiores al 1.1\% en la región de Nueva York.

El efecto de la vegetación como causa dominante de fallas también ha sido ampliamente estudiado. \textcite{ref15} cuantifican el impacto del estándar Enhanced Tree Trimming mediante análisis estadístico y modelos de predicción de interrupciones, mostrando reducciones entre el 16\% y 48\% en fallas durante tormentas. Por su parte, \textcite{ref16} introducen un modelo predictivo para componentes de red ante huracanes basado en regresión logística, mientras que \textcite{ref17} proponen un marco general para extraer conocimiento de datos históricos de la red de Nueva York con el fin de priorizar mantenimiento y estimar la vulnerabilidad de componentes críticos.

Otros aportes relevantes abordan la identificación de causas de fallas. \textcite{ref18} aplican un algoritmo inspirado en sistemas inmunológicos (AIRS) para clasificar causas de interrupción (árboles, animales, rayos), enfrentando el problema de datos desbalanceados. Mientras tanto, \textcite{ref19} establece los fundamentos teóricos de los Random Forest, técnica ampliamente usada en predicción de fallas debido a su robustez y capacidad de manejar alta dimensionalidad.

La estimación de interrupciones bajo condiciones extremas también ha sido examinada mediante enfoques estadísticos. \textcite{ref20} modelan la distribución espacial de fallas durante huracanes en la costa del Golfo utilizando variables físicas medibles. \textcite{ref21} estudian la incertidumbre en modelos de predicción de tormentas y demuestran que el rendimiento de los modelos basados en machine learning depende fuertemente del tamaño y representatividad del conjunto de entrenamiento. A su vez, \textcite{ref22} presentan AdaBoost+, un método de ensamble que supera regresión, redes neuronales y mixture of experts para estimar fallas asociadas al viento y los rayos.

Otras investigaciones incorporan análisis multivariable para predecir tipos específicos de interrupciones. \textcite{ref23} evalúan arquitecturas de aprendizaje automático para predecir fallas por viento, nieve e hielo, mientras que \textcite{ref24} desarrollan un enfoque probabilístico para fallas asociadas a rayos utilizando aprendizaje no supervisado y técnicas de clasificación balanceada.

Asimismo, \textcite{ref25} aplican lógica difusa para análisis predictivo de riesgo, y \textcite{ref26} comparan árboles de regresión no paramétricos (QRF y BART), demostrando ventajas complementarias según la granularidad espacial y las características del evento.

Dentro de los aportes más recientes, \textcite{ref27} desarrollan un enfoque basado en celdas espaciales de $1\,\text{km}^2$ y Random Forest para predecir interrupciones bajo tifones, alcanzando precisiones superiores al 92\%. De forma similar, \textcite{ref28} construyen un modelo de predicción de cantidad de usuarios afectados por tifones, aplicando selección de variables para obtener un modelo simplificado sin sacrificar precisión.

Otras líneas introducen enfoques más estructurales o basados en simulación. \textcite{ref29} proponen un modelo que genera sistemas de distribución sintéticos y estima fallas a nivel de edificio mediante funciones de fragilidad, reduciendo la dependencia en datos privados de operadores. \textcite{ref30} desarrollan un modelo de tormentas de dos etapas para predecir interrupciones en tiempo real, mientras que \textcite{ref31} combinan análisis estadístico y redes neuronales profundas para predecir tiempos de reparación.

\textcite{ref32} demuestran que modelos basados en árboles y ensambles, alimentados con datos atmosféricos generados mediante WRF, pueden predecir con alta resolución espacial la distribución de fallas en Connecticut. Por su parte, \textcite{ref33} incorporan procesamiento de lenguaje natural para actualizar predicciones de duración de interrupciones a partir de reportes de campo.

Más recientemente, \textcite{ref34} sistematiza los principios y particularidades de la aplicación de métodos de aprendizaje automático a la evaluación de la confiabilidad de sistemas de potencia, y desarrolla un prototipo capaz de identificar en tiempo real estados potencialmente peligrosos de la red. Por su parte, \textcite{ref35} presentan un método de cálculo rápido de la confiabilidad de redes de distribución activas basado en tasas de falla dinámicas de los equipos, mientras que \textcite{ref36} integra técnicas híbridas de inteligencia artificial para mejorar la predicción y la gestión de interrupciones, considerando factores meteorológicos, geográficos y operativos. Finalmente, \textcite{ref10} aplican simulación de Monte Carlo para seleccionar la opción óptima de inversión orientada a mejorar el desempeño de confiabilidad en redes de distribución, vinculando así el diagnóstico predictivo con la priorización de inversiones.

Un cuerpo de literatura particularmente relevante para este trabajo adopta explícitamente la discretización del espacio en celdas o cuadrículas de tamaño fijo como unidad de análisis, en lugar de trabajar sobre componentes individuales de la red o sobre agregados administrativos. Esta línea es pionera en la disyuntiva entre resolución espacial y error de predicción, disyuntiva central para la especificación de panel cuadrante $\times$ semana empleada en el presente trabajo.

\textcite{ref68liu} desarrollan uno de los primeros modelos de este tipo, comparando explícitamente un nivel de agregación por código postal (zip code) frente a una malla de celdas de $1\,\text{km}\times 1\,\text{km}$ para predecir interrupciones causadas por huracanes en Carolina del Norte y del Sur; encuentran que, contrario a la intuición de que una mayor resolución mejora la predicción, la agregación a nivel de celda de 1~km introdujo mayores errores que el nivel de código postal, por lo que finalmente seleccionan la especificación menos granular. Este hallazgo temprano anticipa un problema central en la definición de la unidad espacial óptima: celdas demasiado finas pueden diluir la señal disponible en cada unidad y aumentar la dispersión de los datos (muchas más unidades vacías o con pocos eventos), mientras que unidades más gruesas capturan mejor la señal agregada a costa de perder heterogeneidad espacial local.

Como respuesta a esa disyuntiva, un conjunto de estudios posteriores migró hacia celdas de mayor tamaño. \textcite{ref41} desarrollan modelos aditivos generalizados (GAM) sobre una malla de celdas de $2.5\,\text{km}\times 3.7\,\text{km}$ en la costa del Golfo de México, mostrando que las funciones de suavizado no paramétricas de los GAM capturan mejor las relaciones altamente no lineales entre variables meteorológicas (precipitación, humedad del suelo) y el número de interrupciones por celda, en comparación con modelos lineales generalizados (GLM) sobre la misma malla. \textcite{ref37}, trabajando sobre la misma resolución de celda, muestran que modelos de Random Forest —más simples en su formulación que los GAM y GLM previos, y que requieren solo seis variables de entrada (velocidad de ráfaga, duración de vientos fuertes, humedad del suelo, número de clientes y prácticas de poda de árboles)— logran una precisión comparable o superior, evidenciando que, a igualdad de resolución espacial, la elección del algoritmo también condiciona el desempeño predictivo. \textcite{ref42}, por su parte, adoptan una malla más gruesa de $5\,\text{km}\times 5\,\text{km}$ y proponen un modelo de tres etapas: una primera etapa de clasificación binaria para determinar la presencia de interrupciones en cada celda, una etapa intermedia de agrupamiento (clustering) de celdas según la magnitud del evento para abordar la fuerte asimetría de los datos con valores distintos de cero, y una tercera etapa de regresión para predecir el número de interrupciones dentro de cada grupo; esta arquitectura multietapa aborda directamente el problema de desbalance de clases que surge de manera natural en paneles espaciales de alta resolución temporal, donde la mayoría de las celdas-periodo no registran ningún evento de falla —un reto metodológico compartido con el panel cuadrante $\times$ semana del presente trabajo—.

\textcite{ref68mcroberts} extienden esta línea utilizando el trazado censal (census tract) como unidad espacial, e incorporan factores ambientales locales —cobertura vegetal, profundidad de raíces, tipo de suelo— para mejorar modelos de predicción de interrupciones por huracán previamente calibrados a nivel de celda, evidenciando que la elección de la unidad espacial de agregación no solo afecta la resolución del análisis, sino también qué variables explicativas resultan viables de incorporar dado su nivel de disponibilidad geográfica. En una revisión sistemática reciente, \textcite{ref43} sintetizan esta evolución —desde el nivel de código postal y celdas de 1~km, pasando por mallas de $2.5\times 3.7$~km y $5\times 5$~km, hasta el nivel de trazado censal y de ciudad— y muestran, mediante la calibración conjunta de modelos GLM, GAM y Random Forest sobre datos de 1{,}910 ciudades de Texas, Florida, Nueva Jersey y Nueva York, que los modelos existentes tienden a sobrestimar las interrupciones en unidades espaciales de baja población (en el 19.8\% de las ciudades de su muestra, con una razón de sobrestimación promedio de 5.2), y que los Random Forest, si bien evitan estas sobrestimaciones extremas por estar acotados por el rango de los datos de entrenamiento, no logran extrapolar adecuadamente a eventos de viento fuera del rango observado. Este último hallazgo es particularmente relevante para paneles espaciales de alta frecuencia temporal como el aquí propuesto, en la medida en que la capacidad de extrapolación fuera de la muestra —y no solo el ajuste dentro de muestra— condiciona la utilidad del modelo para la priorización de inversión en cuadrantes con historiales de falla poco representados.

En una línea afín mediante representaciones de grafo, \textcite{ref14} y trabajos posteriores sobre redes neuronales de grafos (Graph Neural Networks) han mostrado que incorporar explícitamente las relaciones espaciales de vecindad entre unidades geográficas —y no solo sus atributos individuales— mejora la predicción de interrupciones, particularmente para eventos extremos de baja frecuencia; investigaciones más recientes combinan estas arquitecturas con datos de riesgo arbóreo derivados de sensores LiDAR a nivel de circuito, reforzando la idea de que la dependencia espacial entre unidades vecinas contiene información predictiva adicional a la de los atributos puntuales de cada unidad. De manera complementaria, estudios centrados en la desagregación geográfica fina de los propios indicadores de confiabilidad —examinando la heterogeneidad de SAIDI, SAIFI y CAIDI a nivel de código postal en distintas áreas metropolitanas mediante análisis de clústeres geoespaciales— ilustran cómo la elección de la unidad espacial de agregación condiciona tanto la capacidad de detectar zonas críticas (\emph{hotspots}) de baja confiabilidad como la interpretación de los resultados en términos de priorización de inversión e intervención regulatoria.

En conjunto, esta línea de literatura evidencia una tendencia consistente: el tránsito desde modelos centrados en el componente individual o en agregados puramente administrativos (municipios, circuitos, empresas) hacia esquemas de discretización espacial regular —celdas, cuadrículas o cuadrantes— que permiten capturar de manera más flexible la heterogeneidad espacial de la vulnerabilidad de la red. Esta transición, sin embargo, no está exenta de nuevos retos metodológicos: el tamaño óptimo de la unidad espacial no es neutral frente al desempeño predictivo \parencite{ref68liu}; la dependencia espacial entre celdas vecinas puede aportar información adicional si se modela explícitamente \parencite{ref14}; y el desbalance de clases derivado de la escasez relativa de eventos de falla en paneles de alta granularidad temporal exige arquitecturas de modelación específicas, como los esquemas multietapa \parencite{ref42}. Estos tres retos —selección del tamaño de la unidad espacial, modelación de la dependencia espacial y tratamiento del desbalance de clases en un panel de alta frecuencia— constituyen precisamente los ejes metodológicos centrales que este trabajo aborda al construir un panel espacial de cuadrante $\times$ semana para la predicción de SAIDI/SAIFI en el contexto colombiano.


La revisión anterior permite identificar un conjunto relativamente estable de variables explicativas que reaparecen de manera sistemática en los modelos de predicción de fallas e interrupciones, independientemente de si la variable de respuesta es un conteo de interrupciones, la fracción de clientes afectados o, como en el presente trabajo, indicadores agregados de confiabilidad (SAIDI/SAIFI) a nivel de cuadrante-semana. A continuación se sintetizan estas variables, agrupadas por categoría, junto con la dirección y el mecanismo de la relación que la literatura documenta con la variable pronosticada.

\paragraph{Variables meteorológicas.} La velocidad del viento —típicamente medida como ráfaga máxima de 3~s— es, de manera consistente, la variable individual con mayor poder explicativo en los modelos revisados: \textcite{ref43} la identifican como el predictor más importante en su modelo de Random Forest, por encima de precipitación y densidad poblacional, y su relación con el número de interrupciones es monótonamente creciente y de naturaleza física directa, ya que vientos por encima de aproximadamente 20~m/s superan el umbral de diseño estructural de postes y líneas aéreas \parencite{ref43}. La duración de vientos fuertes por encima de dicho umbral opera de manera similar y suele estar altamente correlacionada con la velocidad máxima, por lo que buena parte de los modelos retienen solo una de las dos variables tras un análisis de colinealidad \parencite{ref43}. La precipitación acumulada (contemporánea y en ventanas previas de 1, 3, 6 o 12 meses, expresada frecuentemente como índice estandarizado de precipitación, SPI) y la humedad del suelo muestran una relación positiva pero \emph{no lineal} con las interrupciones: suelos saturados facilitan el desarraigo de árboles y la caída de líneas ante vientos moderados, mientras que periodos de sequía prolongada debilitan el sistema radicular de los árboles y los vuelve más susceptibles a fallar incluso ante ráfagas menores \parencite{ref41,ref43}; esta no linealidad —efectos marcados tanto en el extremo húmedo como en el extremo seco de la distribución— es precisamente la motivación central para preferir modelos aditivos generalizados (GAM) sobre especificaciones lineales cuando se dispone de variables climáticas de este tipo \parencite{ref41}.

\paragraph{Variables de vegetación y cobertura del suelo.} La proximidad de árboles a las líneas de distribución y el porcentaje de área arbolada dentro de la unidad espacial de análisis mantienen una relación positiva y robusta con la frecuencia de fallas, dado que la caída de árboles y ramas sobre las líneas aéreas es la causa individual más frecuente de interrupciones en redes con alta exposición a vegetación \parencite{ref15,ref43}. En consecuencia, la intensidad y frecuencia de la poda preventiva (\emph{tree trimming}) actúa como variable de exposición inversa: \textcite{ref15} cuantifican que la adopción de estándares de poda mejorada (\emph{Enhanced Tree Trimming}) reduce las tasas de falla asociadas a tormentas entre 16\% y 48\%, por lo que su ausencia o desactualización en los datos de un cuadrante debería interpretarse como un factor que incrementa la probabilidad esperada de interrupción. La profundidad efectiva de raíces también se ha incorporado como proxy de la resistencia estructural de los árboles ante vientos fuertes, con una relación negativa esperada: mayor profundidad radicular implica menor probabilidad de desarraigo \parencite{ref43}.

\paragraph{Variables de infraestructura y exposición.} El número de clientes servidos y, de manera más general, la densidad poblacional de la unidad espacial funcionan como \emph{proxy} de la densidad de componentes de red expuestos (postes, transformadores, tramos de línea), y presentan una relación positiva con el número absoluto de interrupciones, aunque no necesariamente con la tasa o fracción de clientes afectados, distinción que resulta crítica al momento de definir la variable de respuesta del modelo \parencite{ref43}. La topología de la red —redes radiales frente a redes en malla o interconectadas— también se documenta como determinante de la propagación de fallas: sistemas radiales, más comunes en zonas rurales, tienden a mostrar interrupciones de mayor alcance ante la falla de un único componente aguas arriba \parencite{ref43}, un mecanismo que en el caso colombiano resulta pertinente para distinguir cuadrantes urbanos de cuadrantes rurales o de baja densidad.

\paragraph{Variables topográficas y geográficas.} La elevación y variables derivadas de modelos digitales de terreno se han incorporado por su efecto indirecto sobre la velocidad del viento y los patrones de precipitación a escala local \parencite{ref43}, aunque en varios estudios su contribución marginal al poder predictivo resulta baja una vez controlado por viento y precipitación directamente medidos, por lo que su inclusión debe evaluarse en función de la disponibilidad de las variables meteorológicas primarias \parencite{ref43}.

\paragraph{Variables de historial de fallas y autocorrelación.} Finalmente, un conjunto de estudios documenta la relevancia de incorporar la propia historia de interrupciones de la unidad espacial —o de sus vecinas— como predictor autorregresivo. \textcite{ref33} muestran que actualizar predicciones de duración de interrupciones a partir de reportes de campo en tiempo real mejora sustancialmente la precisión respecto a modelos que solo usan información ex ante, mientras que los enfoques de redes neuronales de grafos revisados en la sección anterior \parencite{ref14} evidencian que el número y la severidad de fallas en celdas vecinas aporta información predictiva adicional a la de los atributos propios de cada celda, lo cual sugiere la pertinencia de incluir rezagos espaciales y temporales de SAIDI/SAIFI —y no solo covariables exógenas— en la especificación del panel cuadrante $\times$ semana propuesto en este trabajo.